% §15.2 The BPHF Construction
%
% Cross-reference resolution status:
%
% DEFINED here:
%   \label{sec:bphf-construction}                -- this section
%   \label{def:bphf}                             -- the BPHF definition
%   \label{alg:bphf-build-query}                 -- the build/query pseudocode
%   \label{eq:bphf-r}                            -- rank width
%   \label{eq:bphf-build}                        -- build pass
%   \label{eq:bphf-fnr}                          -- k-perfect FNR law
%
% ALREADY DEFINED (chapter 5):
%   \Cref{def:one-sided}                         -- chapter 5 §5.4
%
% ALREADY DEFINED (chapter 6):
%   \Cref{def:perfect-filter-line}               -- chapter 6 §6.4
%   \Cref{prop:fpr-zero-closure}                 -- chapter 6 §6.4
%
% ALREADY DEFINED (chapter 14):
%   \Cref{def:bhf}                               -- chapter 14 §14.2
%   \Cref{sec:bhf-construction}                  -- chapter 14 §14.2
%   \Cref{alg:bhf-insert-query}                  -- chapter 14 §14.2
%   \Cref{prop:bhf-fpr}                          -- chapter 14 §14.4
%
% ALREADY DEFINED (chapter 15):
%   \Cref{sec:bhf-to-bphf}                       -- §15.1
%
% FORWARD REFS within chapter 15:
%   \Cref{sec:space-time-analysis}               -- §15.3 (Task 4)
%   \Cref{sec:ranked-search}                     -- §15.4 (Task 5)
%
% §15.2 states the BPHF construction in full: two ingredients (perfect
% hash phi and rank table R), build pass, query algorithm with the
% rank-consistency check. It then states honestly what the bare
% construction delivers (it equals the BHF in the general query model),
% identifies the restricted witness model that delivers FPR = 0, and
% introduces k-perfect hashing as the dial that buys a controlled FNR
% on the FPR-zero axis. Worked example at n = 1000.

\section{The BPHF Construction}
\label{sec:bphf-construction}

\Cref{sec:bhf-to-bphf} set the goal: a structure on the FPR-zero axis
of the perfect-filter line, mirroring the way the Bernoulli Hash Filter
sits on the FNR-zero axis. The plan keeps the perfect hash
$\phi : A \to [n]$ from chapter~14 and replaces the fingerprint table
with a \emph{rank table} whose contents the query verifies against the
candidate's own rank. The present section makes the construction
definite, then states plainly what it does and does not deliver. The
honest accounting is the heart of the chapter: the rank table alone
does not flip the error axis, and seeing exactly why is what reveals
the asymmetry between the two halves of the perfect-filter line.

\subsection*{The two ingredients}

Fix a finite universe $\Universe$ and a designated set $A \subseteq
\Universe$ with $|A| = n$. The construction depends on two objects.

\smallskip
\noindent\emph{(i) The perfect hash $\phi$.} The first ingredient is
the same perfect hash function
\[
  \phi \;:\; A \;\longrightarrow\; [n], \qquad [n] = \{1, 2, \ldots, n\},
\]
the chapter~14 construction used: an injection on $A$ that, since
$|A| = n$, bijects the designated set onto its addressing range.
$\phi$ is constructed at build time from $A$ itself and extends in
some canonical way to a function $\phi : \Universe \to [n]$ on the
entire universe. The build cost and storage cost of $\phi$ are
identical to chapter~14's and are computed in
\Cref{sec:space-time-analysis}. For the present section $\phi$ is a
black-box bijection $A \to [n]$ that the construction reuses without
modification.

\smallskip
\noindent\emph{(ii) The rank table $R$ and rank assignment $\rho$.}
The second ingredient is a \emph{rank table}
\[
  R \;:\; [n] \;\longrightarrow\; \{0, 1\}^r,
\]
an array of $n$ slots, each holding an $r$-bit string. The rank
width is
\begin{equation}
\label{eq:bphf-r}
  r \;=\; \lceil \log_2 n \rceil,
\end{equation}
the smallest integer for which $2^r \ge n$. The table is populated
by a \emph{rank assignment}
\[
  \rho \;:\; A \;\longrightarrow\; \{0, 1, \ldots, n-1\},
\]
a function that gives each element of $A$ a distinct integer in
$[0, n)$. The assignment may come from any source the application
supplies: a position in a sorted ordering of $A$, a relevance score
discretized to $n$ buckets, or an external rank derived from
metadata. The distinctness of the ranks is what fixes $r$ at
$\lceil \log_2 n \rceil$: it takes that many bits to label $n$
elements apart.

The two ingredients play complementary roles. $\phi$ is collision-free
on $A$ by construction and addresses the table. $\rho$ is a per-element
value the query can evaluate for any candidate $x$, member or not. The
query compares the stored $R[\phi(x)]$ against the evaluated $\rho(x)$
and reports membership only when the two agree. Where the BHF compared
a stored fingerprint $g(a)$ against the candidate's $g(x)$, the BPHF
compares a stored rank $\rho(a)$ against the candidate's $\rho(x)$. The
substitution is small in the pseudocode. Whether it changes anything in
the error analysis is the question the rest of the section answers, and
the answer is more subtle than the symmetry of the picture suggests.

\subsection*{Construction}

Given $A$, $\phi$, and $\rho$, the BPHF table is built in one pass
over $A$:
\begin{equation}
\label{eq:bphf-build}
  \text{for each } a \in A: \quad R[\phi(a)] \;\coloneqq\; \rho(a).
\end{equation}
The bijectivity of $\phi$ on $A$ means each slot of $R$ is written
exactly once, by exactly the element $a$ for which $\phi(a)$ names
the slot. After the pass, the table $R$ records, for every element
of $A$, the rank $\rho(a)$ stored at the slot $\phi(a)$ reserved for
$a$. No slot remains unwritten because $\phi$ bijects $A$ onto
$[n]$; the table is fully populated.

\begin{definition}[Bernoulli Perfect Hash Filter]
\label{def:bphf}
The \emph{Bernoulli Perfect Hash Filter} (BPHF) for a designated set
$A$ with parameters $\phi$, $\rho$, and $r = \lceil \log_2 n \rceil$
is the data structure $(\phi, R)$ together with the membership query
operation:
\begin{itemize}
  \item \textsc{Query}$(x)$: return $\top$ if $R[\phi(x)] = \rho(x)$,
    return $\bot$ otherwise.
\end{itemize}
After the build pass \eqref{eq:bphf-build}, the predicate
$x \mapsto \textsc{Query}(x)$ defines the membership report of the
approximate set $\ASet{A}$ that the BPHF represents.
\end{definition}

The definition has the same shape as \Cref{def:bhf}: a perfect-hash
addressing function, a per-slot value, and a query that reads one
slot and compares its contents to a value the query computes. The
only structural change is which value is stored. The fingerprint $g$
of chapter~14 was a uniform random $r$-bit string per element. The
rank $\rho$ is a per-element value chosen by the application. The
shapes match; the consequences, as the next subsection shows, do not
match the way the picture promises.

\Cref{alg:bphf-build-query} writes the construction and query in
pseudocode.

\begin{algorithm}[ht]
  \caption{BPHF construction and query}
  \label{alg:bphf-build-query}
  \SetKwInOut{Input}{input}
  \SetKwInOut{Output}{output}
  \SetKwProg{func}{procedure}{}{}
  \SetKwProg{fn}{function}{}{}

  \func{\textsc{Build}($A$, $\rho$)}{
    construct perfect hash $\phi : A \to [n]$\;
    allocate rank table $R$ of size $n$, each slot $r$ bits\;
    \For{each $a \in A$}{
      $R[\phi(a)] \gets \rho(a)$\;
    }
    \Return $(\phi, R)$\;
  }
  \BlankLine
  \fn{\textsc{Query}($x$)}{
    $i \gets \phi(x)$\;
    $v \gets \rho(x)$\;
    \lIf{$R[i] = v$}{\Return $\top$}
    \Return $\bot$\;
  }
\end{algorithm}

\textsc{Build} runs in one pass over $A$ plus the cost of
constructing $\phi$. \textsc{Query} costs two function evaluations
($\phi(x)$ and $\rho(x)$) and an $r$-bit comparison. Against the BHF
pseudocode of \Cref{alg:bhf-insert-query}, the only changed line
substitutes $\rho$ for $g$ in the value being compared. The two
algorithms share their performance profile. The interesting question
is whether they share their error profile, and the answer requires
care.

\subsection*{What the bare construction delivers}

Read the algorithm directly. Take a member $a \in A$. Because $\phi$
bijects $A$ onto $[n]$, the slot $\phi(a)$ was written by exactly one
element, namely $a$ itself, and the build pass stored $R[\phi(a)] =
\rho(a)$. At query time the member recomputes $\phi(a)$ and $\rho(a)$,
reads its own slot, and finds the stored value equal to the recomputed
one. The match always succeeds. A member is never missed:
\[
  \fnrate \;=\; 0.
\]
Two members assigned the same rank do not interfere, because $\phi$
sends them to different slots, and each reads back the rank it wrote.
The rank-collision worry that the symmetry argument might suggest does
not arise: a perfect hash gives every member its own slot.

Now take a non-member $x \notin A$. The extended hash sends $x$ to some
slot $\phi(x) = \phi(a^{*})$ owned by a member $a^{*}$, and that slot
holds $\rho(a^{*})$. The query accepts $x$ exactly when the recomputed
$\rho(x)$ equals the stored $\rho(a^{*})$. If $\rho$ behaves like a
uniform $r$-bit hash on the universe, then $\rho(x)$ is an $r$-bit value
independent of $\rho(a^{*})$, and the two agree with probability
$2^{-r}$. So
\[
  \fprate \;=\; 2^{-r}.
\]
This is exactly the fingerprint-collision probability the BHF analysis
of \Cref{prop:bhf-fpr} bounded. The bare BPHF, queried on the full
universe with a recomputable rank, is the BHF wearing a different
label: zero false negatives, false positives at the floor $2^{-r}$. The
rank table alone does not flip the error axis. Renaming a fingerprint a
rank changes nothing the query can observe.

This is the asymmetry \Cref{sec:bhf-to-bphf} promised, made concrete. A
false positive is a coincidence: one stored value collides with one
recomputed value, and a single $r$-bit slot makes the coincidence rare.
A false-negative-free filter is therefore cheap, and the BHF buys it
for $r$ bits per element. The FPR-zero axis asks for the opposite
guarantee, that no non-member is ever accepted, and that is not a
coincidence to make rare. It is a certificate to enforce against every
one of the $|\Universe| - n$ non-members at once. Certifying a subset
exactly is the expensive direction. The next two subsections show the
two honest routes to it.

\subsection*{How the false-positive rate becomes zero}

The bare construction admits a non-member only when the non-member's
recomputed rank happens to equal the rank stored at the slot it lands
in. Close that gap and the FPR is zero. There is exactly one way to
close it without storing a universe's worth of information: arrange the
query so that a non-member never presents a matching rank in the first
place.

Call $\rho$ a \emph{witness} for $A$ when $\rho$ is defined and
verifiable only for genuine members, in the precise sense that the
recomputed $\rho(x)$ equals the stored rank at slot $\phi(x)$ if and
only if $x$ is the member that claimed that slot. Under a witness, the
acceptance condition $R[\phi(x)] = \rho(x)$ holds only for members, so
\[
  \fprate \;=\; 0
\]
holds exactly, by construction, not in expectation. The relevance
ranking of a search corpus is the standard example, worked in
\Cref{sec:ranked-search}: a query is posed only for a document that
already carries its corpus rank, a document outside the corpus has no
such rank to present, and the witness condition is met.

The witness model is not a free upgrade over the bare construction. It
restricts the query domain to candidates that arrive with a verifiable
rank, and it pays for the FPR-zero guarantee in the only currency that
buys it: the rank table now stores, in its $n$ slots of
$\lceil \log_2 n \rceil$ bits, enough to name every member apart and to
certify any presented witness against the right slot. That is the cost
of certifying a subset exactly, and it is the same counting argument
that set the space lower bounds earlier in the book. There is no
sub-logarithmic shortcut: a structure that accepts no non-member must,
in effect, encode which elements it does accept, and encoding a
designated subset of an arbitrary universe costs on the order of
$\log_2 n$ bits per element once the query domain is the corpus itself.
A genuinely universe-wide FPR-zero filter, one that rejects every one
of the $|\Universe| - n$ outsiders by recomputation alone, would have
to store on the order of $\log_2 |\Universe|$ bits per element, and at
that price it is a dictionary, not a compact filter.

With distinct ranks and a true perfect hash, the witness model places
the BPHF at the \emph{corner} of the perfect-filter line: both error
rates are zero on the query domain, since members always match and
non-members cannot present a witness. The structure is exact there. The
corner is the cleanest point of the FPR-zero axis, and it is also the
least interesting one, because it is not yet an approximate set. The
mirror of the BHF, a tunable structure that trades a controlled error
for space, comes from relaxing the perfect hash.

\subsection*{A controlled false-negative rate via k-perfect hashing}

The BHF moves along its axis by changing the fingerprint width: more
bits per element drive the FPR down toward zero, fewer bits trade
accuracy for space. The BPHF needs the matching dial on its own axis, a
way to spend less space in exchange for a controlled, nonzero FNR while
holding $\fprate = 0$. The dial is the perfect hash itself.

A \emph{$k$-perfect hash} relaxes the injectivity requirement, allowing
up to $k$ elements of $A$ to share a slot. The construction is cheaper
to find and uses fewer slots, but a shared slot can hold only one
member's rank. The members that lose the contention for their slot are
no longer retrievable: each reads a slot holding a different member's
rank, the comparison fails, and they are reported as non-members. These
lost members are the false negatives, and they are the only error the
$k$-perfect BPHF makes. A non-member still cannot present a witness, so
the FPR stays exactly zero.

The false-negative rate is the fraction of members that lose their
slot. With $m$ slots holding $n$ members under a balanced $k$-perfect
construction, roughly $n - m$ members are crowded out, so
\begin{equation}
\label{eq:bphf-fnr}
  \fnrate \;\approx\; 1 - \frac{m}{n},
\end{equation}
which is zero when $m = n$ (the true perfect hash, the exact corner of
the previous subsection) and grows as the table shrinks. A table at
ninety percent of full size, $m = 0.9\,n$, misses about one member in
ten while admitting no non-member at all. This is the mirror of the BHF
made precise: the BHF holds $\fnrate = 0$ and dials $\fprate$ down by
adding fingerprint bits; the BPHF holds $\fprate = 0$ and dials
$\fnrate$ down by adding slots. Each construction lives on its own axis
of the perfect-filter line of \Cref{def:perfect-filter-line}, and each
pays for accuracy in space, on opposite axes, by
\Cref{prop:fpr-zero-closure}.

\subsection*{A worked example}

Fix a designated set of size $n = 1000$. The rank width is
\[
  r \;=\; \lceil \log_2 1000 \rceil \;=\; \lceil 9.966 \rceil \;=\; 10
  \text{ bits},
\]
since $2^{10} = 1024 \ge 1000 > 2^9 = 512$. Ten bits per slot label one
thousand members apart. The rank table occupies
\[
  n \cdot r \;=\; 1000 \cdot 10 \;=\; 10{,}000 \text{ bits} \;=\; 1250
  \text{ bytes},
\]
plus the perfect-hash representation, on the order of $2$ to $3$ bits
per element at $n = 1000$ by the chapter~14 analysis, for a total near
$12{,}000$ to $13{,}000$ bits, about $1500$ to $1625$ bytes.

At the exact corner, a true perfect hash with distinct ranks queried in
the witness model, the error guarantees are
\[
  \fprate \;=\; 0 \quad \text{(by construction, in the witness model),}
  \qquad
  \fnrate \;=\; 0 \quad \text{(every member owns a slot).}
\]
The structure is exact on its query domain: it never accepts a
non-corpus candidate and never misses a corpus member. It is a
perfect-hash dictionary viewed as the corner of the perfect-filter
line.

To turn it into a tunable approximate set on the FPR-zero axis, shrink
the table with a $k$-perfect hash. Take $m = 900$ slots for the same
$n = 1000$ members. By \eqref{eq:bphf-fnr} the false-negative rate is
about
\[
  \fnrate \;\approx\; 1 - \frac{900}{1000} \;=\; 0.10,
\]
so roughly one corpus member in ten is crowded out and missed, while
\[
  \fprate \;=\; 0
\]
holds exactly: no candidate from outside the corpus is ever accepted.
The ten percent smaller table is the price of the ten percent miss
rate, and the affirmative guarantee is untouched.

For comparison, a BHF on the same set at $\fprate = 0.01$ uses $r = 7$
bits per element ($7000$ bits total) and provides $\fnrate = 0$ at
$\fprate \le 0.0078$. The BHF and the BPHF answer opposite needs. The
BHF never misses a member and rarely admits a stranger; the BPHF never
admits a stranger and occasionally misses a member. The choice between
them is the choice of which error an application can tolerate, and the
perfect-filter line says an application that can tolerate neither must
pay the dictionary price for both. The application discussion of
\Cref{sec:ranked-search} works through one setting where the BPHF's
profile is exactly the one the problem wants.
